%% ch02 --- Reguła następnego kroku
%%
%% Napisany we wrześniu 2026 według wzoru rozdziału 1. Każda liczba, której
%% czytelnik nie policzy w pamięci, pochodzi z code/measure/ch02.py, który
%% importuje TE SAME reguły, które drukują listingi (code/ch02/), oraz
%% funkcje `linear` i `orbit` z chaoslab, których te listingi używają.
%%
%% CZEGO TEN ROZDZIAŁ NIE MOŻE POWIEDZIEĆ. Uczy iteracji na regułach
%% LINIOWYCH, których całą przyszłość da się odczytać bez iterowania, i
%% wprowadza nachylenie jako czynnik, przez który mały błąd mnoży się w
%% jednym kroku. Mierzy nachylenie jednej krzywej reguły i na tym poprzestaje:
%% co robi orbita reguły krzywej, która w jednych miejscach rozciąga błędy,
%% a w innych je zmniejsza, zostaje zadane jako pytanie, bez odpowiedzi.
%% Sprzężenie zwrotne i reguła krzywa należą do rozdziału 3, odwzorowanie
%% logistyczne i pajęczyna na krzywej do rozdziału 5, a wykładniczy wzrost
%% błędów do rozdziału 6. Nic w tym rozdziale nie jest chaotyczne.
%%
%% Bliźniak: chapters/en/ch02-iteration.tex. Podrozdział za podrozdziałem,
%% makro za makrem, liczba za liczbą; tools/parity.py je porównuje.

\chapter{Reguła następnego kroku}
\label{ch:iteration}
\index{iteracja}\index{reguła!następnego kroku}

\noindent
Lekarz przepisuje ci jedną tabletkę każdego ranka. Każda tabletka wprowadza
do krwi 100 miligramów leku, a do następnego ranka organizm usuwa połowę
tego, co w niej było. Bierzesz jedną dziennie przez miesiąc, rok, dziesięć
lat: czy ilość leku we krwi rośnie tak długo, jak długo bierzesz tabletki?

Zapamiętaj swoją odpowiedź. To pytanie dotyczy reguły stosowanej raz za
razem \dash{} połowa, tabletka; połowa, tabletka \dash{} a reguły tego
samego rodzaju rządzą kontami oszczędnościowymi, stygnącą herbatą,
populacjami, kredytami i hotelowym prysznicem. Rozdział~\ref{ch:models}
nazwał wiersz \code{temp = next\_minute(temp)} najważniejszym wierszem tej
książki. Ten rozdział mówi o tym, co ten wiersz robi, gdy puścisz go w
nieskończoność, a kończy się tym, że potrafisz powiedzieć, co cała rodzina
reguł zrobi na dłuższą metę, wcale ich nie uruchamiając.

\begin{outcomes}
\outcome{zapisać regułę następnego kroku w notacji tej książki i wyznaczyć
  jej orbitę, ręcznie i w Pythonie}
\outcome{rozpoznać regułę liniową, która stoi za oszczędnościami, lekami i
  stygnięciem, i odczytać jej zachowanie na dłuższą metę z jednej liczby}
\outcome{odpowiedzieć logarytmem na pytanie \enquote{ile kroków, zanim\ldots}}
\outcome{znaleźć stan, w którym reguła stoi w miejscu, i rozstrzygnąć, czy
  zostaje w nim po lekkim pchnięciu}
\outcome{zmierzyć nachylenie reguły w punkcie i odczytać diagram
  pajęczynowy}
\end{outcomes}

\section{Dziś, jutro, pojutrze}
\label{sec:ch02-rule}
\index{stan}\index{orbita}

\noindent
Wpłać 1000 złotych na konto oprocentowane na 5 procent rocznie i nie ruszaj
ich. Po roku bank dopisze 50, czyli 5 procent z 1000, i będziesz mieć 1050.
Po drugim roku dopisze 5 procent od 1050 i będziesz mieć
\val{ch02.save.two}; po trzecim \val{ch02.save.three}. Co roku to samo
dzieje się z tym, co jest na koncie.

Oznacz stan konta po $n$ latach przez $x_n$, co czytamy \enquote{$x$ z
indeksem $n$}: małe $n$ mówi, który to rok, i niczego nie mnoży. Tak więc
$x_0 = 1000$ to punkt wyjścia, $x_1 = 1050$ to stan rok później, a reguła
banku to
\[ x_{n+1} = \num{1.05}\, x_n , \]
czyli stan w przyszłym roku to \num{1.05} razy stan w tym roku, co roku.

Każda reguła krokowa w tej książce ma taki kształt. Jest \emph{stan},
liczba $x$ mówiąca, jak jest teraz; jest \emph{reguła} $f$, która oznacza
\enquote{cokolwiek robisz ze stanem, żeby dostać następny}; i jest
następny stan, czyli $f$ zastosowane do obecnego:
\[ x_{n+1} = f(x_n) . \]
Gdy numer kroku nie ma znaczenia, regułę zapisuje się ze strzałką,
$x \mapsto \num{1.05}\,x$, co czytamy \enquote{$x$ przechodzi w \num{1.05}
razy $x$}. Lista $x_0, x_1, x_2, \ldots$ to \emph{orbita} punktu $x_0$, a
wytwarzanie jej, przez podawanie każdej odpowiedzi z powrotem jako
następnego pytania, to \emph{iteracja}.

\mermaidfig{ch02-feed-back}{Iteracja. Reguła zamienia obecny stan w
następny, który staje się obecnym stanem; na każdym kroku działa ta sama
reguła.}{pętla iteracji}

\begin{think}
Weź regułę $x \mapsto 2x + 1$: podwój, potem dodaj jeden. Zaczynając od
$x_0 = 1$, ile wynoszą $x_1$, $x_2$ i $x_3$?
\end{think}
\reveal{$3$, $7$ i $15$. Dwa razy $1$ plus jeden to $3$; dwa razy $3$ plus
jeden to $7$; dwa razy $7$ plus jeden to $15$.}
Do $x_3$ nie da się dojść inaczej niż przez $x_1$ i $x_2$: sto kroków
naprzód to sto kroków, żmudnych na papierze i bez wysiłku dla maszyny.

\begin{projectbox}
Biblioteka książki, \pkg{chaoslab}, robi te kroki za ciebie. Jej
\code{orbit} to pętla z rozdziału~\ref{ch:models} z rachunkowością
załatwioną raz na zawsze:

\pyregion{code/src/chaoslab/iterate.py}{orbit}{Zastosuj regułę $n$ razy i zachowaj każdy stan}{lst:ch02-orbit}

\noindent
Jej partnerka \code{linear(a, b)} zwraca regułę $x \mapsto a\,x + b$,
gotową do podania funkcji \code{orbit}.
\end{projectbox}

\section{Jedna prosta, trzy losy}
\label{sec:ch02-fates}
\index{reguła liniowa}\index{reguła!liniowa}

\noindent
Reguła oszczędności tylko mnoży. Wiele reguł mnoży, a potem dodaje:

\begin{itemize}
\item \textbf{Herbata} z rozdziału~\ref{ch:models} traci co minutę jedną
  dziesiątą swojej różnicy do pokoju o temperaturze 20 stopni: z $x$
  przechodzi w $x - \num{0.1}(x - 20)$, co po uporządkowaniu daje
  $\num{0.9}\,x + 2$.
\item \textbf{Tabletki}: połowa, potem dodaj 100, czyli
  $x \mapsto \num{0.5}\,x + 100$.
\item \textbf{Populacja}, w której urodzenia co roku przewyższają zgony o
  jedną piątą: $x \mapsto \num{1.2}\,x$.
\item \textbf{Plotka}, którą każdy, kto ją zna, przekazuje dziennie jednej
  nowej osobie: $x \mapsto 2x$.
\end{itemize}

\noindent
Wszystkie mają jeden kształt, $x \mapsto a\,x + b$: pomnóż przez stałą
liczbę $a$, potem dodaj stałą liczbę $b$. Jej wykres, z obecnym stanem na
osi poziomej i następnym na pionowej, jest linią prostą, dlatego nazywa się
ją regułą \emph{liniową} \dash{} choć, jak zaraz zobaczysz, to, co robi w
czasie, jest przeważnie krzywe. Uruchom cztery takie reguły z tego samego
miejsca:

\pyregion{code/ch02/fates.py}{rules}{Cztery reguły liniowe iterowane obok siebie z tego samego startu}{lst:ch02-fates}

\transcript{ch02-fates}

\noindent
Pierwsza kolumna podwaja się w nieskończoność, a druga maleje o połowę ku
zeru. Trzecia skrada się do $10$ od dołu. Czwarta też ustala się na $10$,
ale skacząc z jednej strony na drugą, a każdy skok jest mniejszy od
poprzedniego: to trzecia kolumna w lustrzanym odbiciu. Oto trzy losy reguły
liniowej \dash{} rosnąć bez ograniczeń, zanikać do zera albo się ustalić
\dash{} a ujemny mnożnik dodaje naprzemienność, ale nie czwarty los.

\plotfig{ch02-three-fates}{Trzy losy reguły liniowej przy tym samym
starcie. Jedna rośnie bez ograniczeń, jedna zanika, a dwie ustalają się na
tej samej liczbie: jedna się do niej skrada, druga skacze z boku na
bok.}{trzy reguły liniowe i ich trzy losy}

\begin{think}
Wróćmy do konta oszczędnościowego. Chcesz znać stan po dziesięciu latach,
nie robiąc dziesięciu kroków. Jakie jedno obliczenie go daje?
\end{think}
\reveal{Pomnóż 1000 przez \num{1.05} dziesięć razy z rzędu: $x_{10} =
\num{1.05}^{10} \times 1000$, co daje \val{ch02.save.ten}.}
Małe uniesione $10$ to \emph{potęga}: $\num{1.05}^{10}$ to \num{1.05}
pomnożone przez siebie dziesięć razy. Gdy reguła tylko mnoży, $x \mapsto
a\,x$, jej orbita ma wzór, $x_n = a^n x_0$, i można przeskoczyć do
dowolnego kroku. Wzrost, w którym każdy krok mnoży przez tę samą liczbę,
nazywa się \emph{wykładniczym}; jego malejący bliźniak, z $a$ między $0$ a
$1$, to \emph{zanik wykładniczy}.

\section{Powolne eksplozje i liczenie kroków}
\label{sec:ch02-until}
\index{wzrost wykładniczy}\index{czas podwojenia}\index{logarytm}
\index{okres półtrwania}

\noindent
Słowo \enquote{wykładniczy} stało się synonimem słowa \enquote{szybki}. Oto
test tego przekonania.

\begin{think}
Dwa konta oszczędnościowe zaczynają od 1000. Konto A co roku dopisuje 50.
Konto B rośnie o 2 procent rocznie, czyli wykładniczo. Które jest do przodu
po trzydziestu latach?
\end{think}
\reveal{Konto A, z $1000 + 30 \times 50 = 2500$ wobec
\val{ch02.acc.grows.thirty} na koncie B. B wyprzedza A dopiero w roku
\val{ch02.acc.cross}.}
Jeśli wybrałeś B, rozumowałeś na podstawie słowa, a nie reguły: 2 procent z
1000 to 20, mniej niż 50 konta A, a B potrzebuje dziesięcioleci wzrostu na
własnym wzroście, żeby je dogonić.

\begin{trapbox}
\textbf{\enquote{Wzrost wykładniczy to szybki wzrost.}} To wzrost
proporcjonalny do tego, co już jest, a mnożnik może być bardzo bliski
jedności. Przy 1 procencie rocznie suma podwaja się po mniej więcej
\val{ch02.dbl.one.log} latach. Wzrost wykładniczy gwarantuje nie szybkość,
tylko nieuchronność: jakkolwiek mała byłaby stopa, prześcignie każdy
wzrost, który dodaje stałą kwotę \dash{} w końcu, a to \enquote{w końcu}
może trwać całe życie.
\end{trapbox}

\plotfig{ch02-two-accounts}{Dodawanie stałej kwoty wobec wzrostu o stały
ułamek. Konto wykładnicze jest z tyłu przez większą część stulecia, a potem
już nie.}{wzrost liniowy wobec wykładniczego}

\noindent
Pożyteczne pytanie o wzrost wykładniczy brzmi \enquote{ile kroków, zanim
się podwoi?}, a odpowiedź nie zależy od punktu startu.

\begin{think}
Przy 7 procentach rocznie po mniej więcej ilu latach suma się podwaja: po
około 5, około 10 czy około 14?
\end{think}
\reveal{Po około 10. Licząc pełne lata, 1000 po raz pierwszy przekracza
2000 w roku \val{ch02.dbl.seven.count}; dokładna odpowiedź to
\val{ch02.dbl.seven.log} roku.}
Jest na to przycisk. Pytanie brzmi: dla jakiego $n$ zachodzi
$\num{1.07}^n = 2$? Na dokładnie takie pytania odpowiada \emph{logarytm}.
Kalkulator ma przycisk oznaczony \code{ln}, logarytm naturalny, a jedyne,
co musisz o nim wiedzieć, to co robi z potęgą:
\[ \ln (a^n) = n \ln a . \]
Zamienia \enquote{pomnóż przez $a$, $n$ razy} na \enquote{dodaj $\ln a$, $n$
razy}. Zlogarytmuj obie strony równania $a^n = 2$, a dostaniesz
$n \ln a = \ln 2$, czyli
\[ n = \frac{\ln 2}{\ln a} : \]
liczba kroków to logarytm tego, jak daleko chcesz dojść, podzielony przez
logarytm jednego kroku. Skoro $\ln 2 = \val{ch02.ln.two}$, a
$\ln \num{1.07} = \val{ch02.ln.seven}$, wychodzi \val{ch02.dbl.seven.log}.

\pyregion{code/ch02/doubling.py}{count}{Po ilu latach stan konta się podwoi: liczone i z logarytmu}{lst:ch02-doubling}

\transcript{ch02-doubling}

\noindent
Liczba kroków to logarytm zaokrąglony w górę. Przy 2 procentach to
\val{ch02.dbl.two.log} wobec \val{ch02.dbl.two.count}: po 35 latach na
koncie jest \val{ch02.dbl.two.short} raza tyle, ile wpłacono, o włos mniej
niż podwojenie. Ostatnia kolumna, \emph{reguła 70}, dzieli 70 przez liczbę
procent; działa, bo logarytm jednego kroku jest bliski samej małej stopie
($\ln \num{1.07}$ to prawie \num{0.07}), a $100 \times \ln 2$ to mniej
więcej 70.

Z maleniem jest tak samo. Lek, z którego organizm co godzinę usuwa jedną
piątą, zachowuje co godzinę \num{0.8}, a czas, po którym zostaje połowa,
jego \emph{okres półtrwania}, wynosi $\ln \num{0.5} / \ln \num{0.8}$, czyli
około \val{ch02.half.fifth} godziny.

\begin{lab}{l02_01_until}{Ile kroków do celu?}
Napisz \code{steps\_until}, która liczy kroki reguły $x \mapsto a\,x$ od
$1$, aż $x$ osiągnie cel, oraz \code{time\_until}, która zadaje to samo
pytanie logarytmowi, dla wzrostu i dla zaniku. Potem napisz
\code{half\_life}. Logarytm naturalny w Pythonie to \code{math.log}.
\end{lab}

\section{Gdzie reguła stoi w miejscu}
\label{sec:ch02-fixed}
\index{punkt stały}

\noindent
Teraz pytanie z początku rozdziału. Stanem jest ilość leku we krwi zaraz po
porannej tabletce, w miligramach: $x \mapsto \num{0.5}\,x + 100$, od
$x_0 = 100$.

\begin{think}
Policz kilka pierwszych poranków: $100$, potem połowa tego plus $100$ i tak
dalej. Czy ilość leku rośnie tak długo, jak długo bierzesz tabletki?
\end{think}
\reveal{Nie. Idzie $100$, $150$, $175$, \num{187.5}, \num{193.75} i dalej,
zbliżając się do $200$, ale nigdy go nie osiągając.}
Dlaczego $200$? Bo reguła zostawia $200$ bez zmian: połowa z $200$ to
$100$, a tabletka od razu dokłada $100$ z powrotem. Stan, którego reguła nie
zmienia, to \emph{punkt stały}, a znajdujesz go, pytając, które $x$ spełnia
\[ x = f(x) . \]
Dla tabletek to $x = \num{0.5}\,x + 100$; odejmij $\num{0.5}\,x$ od obu
stron, a dostaniesz $\num{0.5}\,x = 100$, czyli $x = 200$. Dla każdej
reguły liniowej te same kroki dają
\[ x^{*} = \frac{b}{1 - a} , \]
gdzie gwiazdka oznacza punkt stały. (Jeśli $a$ wynosi dokładnie $1$,
reguła tylko dodaje i nigdzie nie stoi w miejscu.) Punkt stały reguły
herbaty leży w $2 / \num{0.1} = 20$ stopniach: w temperaturze pokoju.

\pyregion{code/ch02/drug.py}{tablets}{Jedna tabletka co rano i odległość od punktu stałego}{lst:ch02-drug}

\transcript{ch02-drug}

\noindent
Niedobór w prawej kolumnie co rano maleje o połowę, bo organizm usuwa
połowę tego, co jest, a tabletka zawsze dokłada tych samych $100$. Po
\val{ch02.drug.within} dniach od $200$ dzieli cię mniej niż jeden miligram,
a dokładnie tam nie dotrzesz nigdy. Ilość leku nie rośnie w nieskończoność;
wyrównuje się w punkcie stałym.

\begin{worldbox}
Farmakolodzy nazywają ten punkt stały \emph{stanem stacjonarnym}, a
praktyczna reguła, której uczy się lekarzy i farmaceutów, mówi, że lek
przyjmowany w regularnych odstępach osiąga go po czterech lub pięciu
okresach półtrwania. To tabela powyżej: po czterech okresach półtrwania
niedobór zmalał o połowę cztery razy, co daje poziom \val{ch02.drug.four}
procent stanu stacjonarnego, a po pięciu \val{ch02.drug.five}. Gdy lek jest
potrzebny od razu w pełnej sile, lekarz może podać większą pierwszą dawkę,
\emph{dawkę nasycającą}, która zaczyna orbitę w punkcie stałym, a nie
poniżej niego.
\end{worldbox}

\section{Zostać na miejscu albo uciec}
\label{sec:ch02-stability}
\index{punkt stały!stabilny}\index{punkt stały!niestabilny}\index{stabilność}

\noindent
Oto punkt stały innego charakteru. Masz \num{10000} długu oprocentowanego
na 1 procent miesięcznie i spłacasz 100 miesięcznie, więc co miesiąc dług
przechodzi w $\num{1.01}\,x - 100$.

\begin{think}
Co dzieje się z długiem przez lata? A co, jeśli jesteś winien \num{11000},
przy tym samym oprocentowaniu i tej samej racie?
\end{think}
\reveal{Przy \num{10000} nic, nigdy: odsetki wynoszą dokładnie 100, a twoje
100 je znosi. Przy \num{11000} odsetki wynoszą 110, twoje 100 ich nie
pokrywa i po dziesięciu latach jesteś winien \val{ch02.loan.up}.}
Punktem stałym tej reguły jest \num{10000}, ale reguła wcale cię do niego
nie przyciąga. Zacznij tysiąc powyżej, a dług ucieka w górę; zacznij tysiąc
poniżej, od \num{9000}, a ucieka w dół, do zera, spłacony po
\val{ch02.loan.payoff} miesiącach, czyli po mniej więcej
\val{ch02.loan.years} latach.

\begin{trapbox}
\textbf{\enquote{Punkt stały to miejsce, w którym układ w końcu się
znajdzie.}} To miejsce, w którym układ pozostaje, jeśli jest w nim
dokładnie. Czy układ w pobliżu zbliża się do niego, czy się oddala, to
osobne pytanie, a punkt stały kredytu odpycha każdy inny stan długu. Punkt
stały, który przyciąga pobliskie stany, jest \emph{stabilny}; taki, który
je odpycha, jest \emph{niestabilny}. Znalezienie punktu stałego mówi ci,
gdzie szukać; dopiero jego stabilność mówi, czy cokolwiek tam będzie.
\end{trapbox}

\noindent
Różnicę robi jedna liczba. Prześledź przez jeden krok odległość między
stanem $x$ a punktem stałym $x^{*}$:
\[ f(x) - x^{*} = (a\,x + b) - (a\,x^{*} + b) = a\,(x - x^{*}) . \]
Pierwszy znak równości korzysta z tego, że $x^{*}$ jest stały, więc
$a\,x^{*} + b$ równa się $x^{*}$. Wyraz $b$ się skraca, a to, co zostaje,
mówi: \textbf{każdy krok mnoży odległość przez $a$}. Tabletki ją połowią;
kredyt powiększa ją o 1 procent miesięcznie; herbata, z $a = \num{0.9}$,
to powód, dla którego dwa kubki z rozdziału~\ref{ch:models} zbliżały się
do siebie co minutę o jedną dziesiątą.

Liczy się wielkość $a$, zapisywana $\abs{a}$, która pomija znak:
$\abs{\num{-0.5}} = \num{0.5}$. Poniżej 1 odległość maleje i punkt stały
jest stabilny; powyżej 1 rośnie i punkt stały jest niestabilny. Znak mówi
tylko, czy odległość zostaje po tej samej stronie, czy przeskakuje na
drugą.

Ujemne $a$ czułeś pod hotelowym prysznicem. Chcesz 38 stopni, woda ma 30, a
ty przesadzasz: każdy obrót kurka zmienia temperaturę o półtora raza tyle,
ile wynosi błąd, w przeciwną stronę. Reguła to $x \mapsto x - \num{1.5}(x -
38)$, czyli $x \mapsto \num{-0.5}\,x + 57$. Z 30 przechodzisz do 42, 36, 39,
\num{37.5}: za gorąco, za zimno, za gorąco, a każde pudło jest o połowę
mniejsze od poprzedniego. Szarpnij kurkiem o dwa i pół raza tyle, ile
wynosi błąd, a $a = \num{-1.5}$: 30, 50, 20, 65, za każdym razem gorzej.

\begin{center}
\begin{tabularx}{\linewidth}{@{}lXX@{}}
\toprule
mnożnik $a$ & pchnięcie od $x^{*}$ & orbita \\
\midrule
powyżej $1$ & rośnie & ucieka, po jednej stronie \\
między $0$ a $1$ & maleje & zbliża się, z jednej strony \\
między $-1$ a $0$ & maleje, zmieniając stronę & ustala się, przeskakując \\
poniżej $-1$ & rośnie, zmieniając stronę & ucieka, w coraz szerszych
wahaniach \\
\bottomrule
\end{tabularx}
\end{center}

\begin{lab}{l02_02_fixed}{Gdzie się zatrzyma i czy tam zostanie?}
Dla reguły $x \mapsto a\,x + b$ napisz \code{fixed\_point};
\code{is\_stable}, która rozstrzyga, czy pchnięcie od punktu stałego maleje
(spróbuj: zacznij trochę obok, zrób jeden krok, porównaj odległości); oraz
\code{steps\_to\_settle}.
\end{lab}

\section{Nachylenie: ile razy mnoży się pchnięcie}
\label{sec:ch02-slope}
\index{nachylenie}

\noindent
Mnożnik $a$ ma drugą nazwę, której używa reszta książki. Narysuj regułę
tabletek jako wykres, z obecnym stanem na osi poziomej i następnym na
pionowej: to prosta $y = \num{0.5}\,x + 100$. Przesuń się wzdłuż osi
poziomej ze 150 do 151, o jedną jednostkę, a prosta wzniesie się ze 175 do
\num{175.5}, o pół jednostki. \emph{Nachylenie} wykresu w punkcie to
właśnie to: \textbf{o ile przesuwa się wynik na każdą jednostkę
przesunięcia wejścia}. Prosta ma wszędzie to samo nachylenie, a dla
$x \mapsto a\,x + b$ to nachylenie wynosi $a$. Zatem odkrycie z
poprzedniego podrozdziału, w słowach, których używa reszta książki,
brzmi: \emph{mały błąd jest na każdym kroku mnożony przez nachylenie}.

Krzywa też ma nachylenie, ale zmienia się ono od punktu do punktu, tak jak
zbocze jest w jednym miejscu strome, a w innym prawie płaskie. Mierzysz je
w punkcie, pchając: przesuń wejście o maleńką wielkość $h$ w obie strony od
$x$, zobacz, o ile przesunął się wynik, i podziel,
\[ \text{nachylenie w } x \approx \frac{f(x + h) - f(x - h)}{2h} , \]
gdzie $\approx$ znaczy \enquote{prawie równe}. W funkcji \code{slope} z
\pkg{chaoslab} pchnięcie $h$ to jedna milionowa: przyrost w górę podzielony
przez przyrost w bok, na odcinku zbyt krótkim, by krzywa zdążyła się
wygiąć.

\pyregion{code/ch02/slope.py}{measure}{Nachylenie zmierzone pchnięciem, na prostej i na krzywej}{lst:ch02-slope}

\transcript{ch02-slope}

\noindent
Nachylenie prostej wynosi \num{0.5}, gdziekolwiek je zmierzysz; krzywa
reguła $x \mapsto x \times x$ jest tym bardziej stroma, im dalej w prawo.

\begin{think}
Reguła kwadratu ma nachylenie $3$ w punkcie \num{1.5}. Chciałeś podać jej
\num{1.5}, a podałeś \num{1.501}, z błędem \num{0.001}. O ile mniej więcej
pomyli się wynik?
\end{think}
\reveal{O mniej więcej \num{0.003}: błąd mnoży się przez nachylenie.
Dokładnie: $\num{1.501} \times \num{1.501} - \num{2.25} =
\val{ch02.nudge.exact}$.}
To fakt, który stoi za każdą odległością w tym rozdziale, teraz na
krzywej: błąd mnoży się przez nachylenie w miejscu, w którym akurat jest.
Na prostej ten czynnik się nie zmienia i dlatego ten rozdział może wszystko
przewidzieć. W regule kwadratu się zmienia: między $0$ a \num{0.5}
nachylenie jest mniejsze niż 1 i błędy maleją, powyżej \num{0.5} jest
większe niż 1 i rosną. Oto więc pytanie, na które ten rozdział nie odpowie.
Co robi orbita, jeśli jej reguła jest krzywa, tak że orbita wciąż
przechodzi z miejsc, które zmniejszają błędy, do miejsc, które je
rozciągają? Rozdział~\ref{ch:feedback} wygina prostą, a
rozdział~\ref{ch:logistic} śledzi orbitę, która robi dokładnie to.

\section{Zobaczyć to: diagram pajęczynowy}
\label{sec:ch02-cobweb}
\index{diagram pajęczynowy}

\noindent
Iterację można oglądać na wykresie samej reguły, z ołówkiem i bez
rachunków. Narysuj prostą reguły i przekątną $y = x$, której wysokość
zawsze równa się położeniu na osi poziomej. Potem:

\begin{enumerate}
\item Od $x_0$ na osi poziomej idź w górę lub w dół do prostej. Wysokość,
  na którą dojdziesz, to $x_1$.
\item Idź poziomo do przekątnej, gdzie wysokość równa się położeniu:
  stoisz teraz nad $x_1$, gotowym do użycia jako następne wejście.
\item W górę lub w dół do prostej daje $x_2$; w bok do przekątnej i tak
  dalej.
\end{enumerate}

\noindent
Narożniki tej ścieżki to narysowana orbita. Funkcja \code{cobweb} z
\pkg{chaoslab} je wypisuje; dla tabletek, zaczynając od pustego krwiobiegu:

\pyregion{code/ch02/cobweb_steps.py}{corners}{Narożniki pajęczyny dla tabletek}{lst:ch02-cobweb}

\transcript{ch02-cobweb}

\noindent
Narożniki na przekątnej niosą orbitę, każdą liczbę dwa razy: $0$, $100$,
$150$, $175$. Ścieżka może się zatrzymać tylko tam, gdzie prosta przecina
przekątną, bo tam $x = f(x)$: przecięcie to punkt stały, który widać.

\begin{think}
Prosta kredytu, $y = \num{1.01}\,x - 100$, jest bardziej stroma niż
przekątna. Czy ścieżka, która zaczyna trochę z boku przecięcia, wspina się
ku przecięciu, czy się od niego oddala?
\end{think}
\reveal{Oddala się. Każdy krok mnoży odległość od przecięcia przez
\num{1.01}, więc ścieżka ląduje za każdym razem trochę dalej, po tej
stronie, po której zaczęła.}

\plotfig{ch02-cobweb}{Pajęczyny na trzech prostych, które przecinają
przekątną w $200$. Nachylenie o wielkości mniejszej niż jeden wciąga
ścieżkę do przecięcia \dash{} jako schodki, gdy nachylenie jest dodatnie,
jako pajęczynę, gdy ujemne \dash{} a nachylenie większe niż jeden ją
wypycha.}{pajęczyna na prostej}

\noindent
Gdy nachylenie jest ujemne, ścieżka owija się wokół przecięcia i to od tej
pajęczyny diagram wziął nazwę. Trzy panele to narysowana tabela z
podrozdziału~\ref{sec:ch02-stability}.

Masz więc coś rzadkiego w tej książce: pełną zdolność przewidywania. Mając
$x \mapsto a\,x + b$ i punkt startu, możesz bez robienia choćby jednego
kroku powiedzieć, co reguła będzie robić zawsze \dash{} ustali się w
$b/(1 - a)$, jeśli $\abs{a} < 1$, ucieknie, jeśli $\abs{a} > 1$, będzie się
wahać naprzemiennie, jeśli $a$ jest ujemne \dash{} i jak długo to potrwa.
Tę pewność odbiera reszta książki: wygnij prostą, a reguła długości jednej
linijki potrafi robić rzeczy, których nie przewidzi żaden wzór z tego
rozdziału.

\begin{summarybox}
\begin{itemize}
\item \textbf{Reguła następnego kroku} zamienia obecny stan w następny,
  $x_{n+1} = f(x_n)$. Powtarzanie jej to \textbf{iteracja}; lista stanów to
  \textbf{orbita}.
\item Oszczędności, stygnięcie i dawkowanie leków podlegają
  \textbf{regule liniowej}, $x \mapsto a\,x + b$: rośnie bez ograniczeń,
  zanika do zera albo się ustala; ujemne $a$ sprawia, że
  \textbf{oscyluje naprzemiennie}.
\item Wzrost \textbf{wykładniczy} jest proporcjonalny, niekoniecznie
  szybki. \textbf{Logarytm} mówi, ile kroków zajmuje zmiana o czynnik $r$,
  $\ln r / \ln a$: stąd czas podwojenia i okres półtrwania.
\item \textbf{Punkt stały} to stan, którego reguła nie zmienia: rozwiąż
  $x = f(x)$, co dla reguły liniowej daje $b/(1 - a)$.
\item Każdy krok mnoży odległość od punktu stałego przez $a$: punkt jest
  \textbf{stabilny}, gdy $\abs{a} < 1$, a \textbf{niestabilny}, gdy
  $\abs{a} > 1$, i w niestabilnym nic nie kończy.
\item \textbf{Nachylenie}, czyli o ile przesuwa się wynik na jednostkę
  przesunięcia wejścia, to czynnik, przez który mały błąd mnoży się w
  jednym kroku. Nachylenie krzywej zmienia się od punktu do punktu.
\item \textbf{Diagram pajęczynowy} rysuje orbitę na wykresie reguły; na
  prostej zbliża się do przecięcia z przekątną, gdy nachylenie ma wielkość
  mniejszą niż jeden.
\end{itemize}
\end{summarybox}

\canyou

\begin{checkpoint}
\item Dla reguły $x \mapsto 3x - 2$ i startu $x_0 = 2$ ile wynoszą $x_1$,
  $x_2$ i $x_3$?
  \answerto{$4$, $10$ i $28$.}
\item Herbata traci co minutę jedną piątą różnicy do pokoju o temperaturze
  20 stopni. Zapisz jej regułę w postaci $x \mapsto a\,x + b$.
  \answerto{$x \mapsto x - \num{0.2}(x - 20)$, czyli $\num{0.8}\,x + 4$.}
\item Miasto rośnie o 3 procent rocznie. Oszacuj czas podwojenia regułą
  70, a potem powiedz, co daje liczenie pełnych lat.
  \answerto{Około $70 / 3$, czyli 23 lata. Licząc lata, liczba mieszkańców
  po raz pierwszy przekracza podwojenie w roku \val{ch02.dbl.three.count};
  logarytm daje \val{ch02.dbl.three.log}.}
\item Znajdź punkt stały reguły $x \mapsto \num{0.8}\,x + 4$. Czy jest
  stabilny?
  \answerto{$\num{0.2}\,x = 4$, więc $x = 20$, temperatura pokoju;
  stabilny, bo $\abs{\num{0.8}} < 1$.}
\item Reguła $x \mapsto 3x - 2$ ma punkt stały w $1$. Czy jest stabilny? Co
  pokazuje pierwsze pytanie?
  \answerto{Niestabilny: każdy krok potraja odległość. Orbita z $2$ jest o
  jeden od punktu stałego, potem o $3$, $9$ i $27$: $4$, $10$, $28$.}
\item Reguła prysznica $x \mapsto \num{-0.5}\,x + 57$ ma nachylenie
  \num{-0.5}. Jest ci o 2 stopnie za zimno. Gdzie jesteś po jednym obrocie
  kurka?
  \answerto{1 stopień za gorąco: błąd zmalał o połowę i zmienił stronę.}
\end{checkpoint}
